\tutorial{3}{The soil-zone coupling}

\begin{goal}
You let the water table of Tutorial Creek rise into Raven's soil (\cmd{:SoilZone PHYSICAL}), see what changes against the
default \cmd{CONCEPTUAL} coupling, add unsaturated columns between the soil and a deep water table
(\cmd{:SoilZoneVadose RICHARDS}, \cmd{KINEMATIC}, \cmd{LAG}) and depth classes, and learn what to look at in
\cmd{GWSoilExchange.csv}, \cmd{GWBudget.csv} and the plots of \cmd{plot\_soilzone.py}. Time: about 45 minutes, of which
about \V{t3:secclasses} seconds for the longest run.
\end{goal}

\section{Two ways to couple the soil and the aquifer}
\begin{description}
\item[\cmd{CONCEPTUAL} (default, Tutorial~2)] The soil and the aquifer meet only through recharge (whatever Raven puts into
  \cmd{GROUNDWATER}) and a seepage drain at the land surface: where the water table reaches the surface, MODFLOW discharges
  groundwater back to Raven. The water table never enters the soil.
\item[\cmd{PHYSICAL}] Raven owns the listed soil layers; MODFLOW's model top is the base of those layers (land surface minus
  their thickness). The water table can rise into the soil: there recharge is rejected (it stays in the soil), saturated
  soil drains to the aquifer, groundwater rises into the soil, and a full soil zone overflows to the surface. Within each
  day the soil zone and MODFLOW are iterated until they agree. On sloping or hummocky land each cell samples the DEM at
  5$\times$5 points, so the low points saturate first.
\end{description}
Below the soil zone, recharge can cross an unsaturated zone: \cmd{:SoilZoneVadose} adds a column per cell between the
soil-zone base and a water table deeper than 5~cm, which delays and spreads the recharge (Richards' equation, a kinematic
wave, or a linear storage cascade).

\section{Start}
Tutorial~3 starts from the finished model of Tutorial~2:
\begin{shell}
$ cp -r $TUT/t3_soil_zone/start ~/work/t3
$ cd ~/work/t3
\end{shell}
First run it as it is (\cmd{CONCEPTUAL}), into its own output folder, to compare with later:
\term{t3_run_conceptual.txt}

\section{The soil zone (\texttt{:SoilZone PHYSICAL})}
\subsection{Soil parameters}
The soil zone will be \cmd{TOPSOIL} and \cmd{SUBSOIL} (\cmd{SOIL[0]} and \cmd{SOIL[1]}): recharge leaves from the bottom layer,
\cmd{SUBSOIL}, by a vertical process (\cmd{:Percolation \dots SUBSOIL GROUNDWATER}), as the coupling requires. Each soil-zone
layer needs \cmd{POROSITY} and \cmd{FIELD\_CAPACITY}; the class of the bottom layer needs \cmd{HYDRAUL\_COND} [mm/d], which with
\cmd{K\_VERT} of the top aquifer layer sets the conductance between soil and aquifer; capillary rise (on by default with
unsaturated columns, Section~\ref{sec:t3vz}) uses \cmd{CLAPP\_B} and \cmd{AIR\_ENTRY\_PRESSURE} [mm]. Add the three columns
to the soil parameter table of \cmd{Creek.rvp}:
\diff{t3_soil.diff}

\subsection{Switch it on}
Add one command at the end of \cmd{Creek.rvg}:
\diff{t3_physical.diff}
and run into a new folder:
\term{t3_run_physical.txt}
The run takes about \V{t3:secphysical}~seconds instead of \V{t3:secconceptual}: every day the soil zone and MODFLOW are
iterated to agreement.
\begin{trouble}
\begin{itemize}
\item \emph{``soil class \dots of SOIL[1] \dots needs HYDRAUL\_COND''}: the \cmd{.rvp} edit is missing.
\item \emph{``RECHARGE\_DELAY \dots cannot be used with :SoilZone PHYSICAL''}: a profile has a recharge delay; set it to 0 (use
      \cmd{:SoilZoneVadose LAG} instead).
\item \emph{``soil-zone interface did not converge''} (rare): raise \cmd{:SoilZoneIteration}, increase
      \cmd{SEEPAGE\_SMOOTHING\_DEPTH}, or shorten the time step.
\end{itemize}
\end{trouble}

\subsection{Read the summary and the warnings}
\term{t3_summary.txt}
\V{t3:nszhru} HRUs have a soil zone; the lake HRU has no soil and keeps the seepage drain. The soil zone is 1.3~m thick under the
valley and the wetland and 1.0~m under the slopes (the soil profiles of the \cmd{.rvp}), and Raven keeps one soil-zone
column for every link between an HRU and a cell. The \V{t3:nwarn} warnings are worth reading:
\term{t3_warnings.txt}
\begin{itemize}
\item \emph{The soil-zone base lies up to \V{t3:offmax}~m from the model top on \V{t3:offlinks} links}: a cell's top is
      the land surface minus the overlap-weighted soil thickness of all the HRUs in it. Cells shared by the valley
      (1.3~m) and a slope (1.0~m), or by an HRU with soil and the lake (none), have a top between the two; the warning says
      where. With HRUs whose soil profiles differ, this is expected.
\item \emph{The specific yield differs from the drainable porosity of the bottom soil layer on \V{t3:sylinks} links}
      (ratio down to \V{t3:syratio}): the slopes' till has $S_y=0.08$, the subsoil a drainable porosity of
      $0.40\times(1-0.25)=0.30$. A water table rising from the till into the soil would change speed there. On these slopes the
      water table stays far below the soil, so it does not matter here; on wet ground make the two agree.
\end{itemize}

\subsection{What changed}
\cmd{GWSoilExchange.csv} has one line per HRU and day:
\term{t3_soilex_head.txt}
The soil-zone base, the soil water level and the water table (means over the HRU's columns), the saturated fraction (the
share of the HRU whose soil the water table reaches) and the exchange in both directions. The package's plotting tool
draws them:
\term{t3_plot_soilzone.txt}
\fig{t3_soilzone.png}{\cmd{plot\_soilzone.py} for the valley HRUs 6 and 1. Left: the water table rises above the soil-zone
base every winter and the soil water level follows it. Right: the exchange (aquifer to soil, positive) and the saturated
fraction (dotted). Below: the totals of all HRUs, rejected recharge and overflow, and the interface iteration.}
\begin{tocheck}
\begin{itemize}
\item Once the water table is in the soil, the soil water level should follow it (left panels).
\item The exchange should change sign only at physical turning points, not from day to day.
\item The saturated fraction of HRU~6 rises and falls with the seasons (between \V{t3:F6min} and \V{t3:F6max}, mean
      \V{t3:F6mean}): with 25 DEM points per cell, the hollows of the valley floor saturate first.
\item Interface updates per step stay small and the last change of the soil level below the tolerance (bottom panels).
\end{itemize}
\end{tocheck}

\section{Unsaturated columns (\texttt{:SoilZoneVadose})}\label{sec:t3vz}
Under the slopes the water table lies far below the soil (a mean depth of \V{t3:physical:wt7}~m under HRU~7). With \cmd{:SoilZoneVadose NONE} (the default) recharge
leaving the soil reaches it the same day. An unsaturated column per cell carries it down instead. The columns need the
retention curve of the aquifer materials: \cmd{POROSITY} (already there), \cmd{CLAPP\_B} and \cmd{AIR\_ENTRY\_PRESSURE} [m]:
\diff{t3_aquifer.diff}
The till values give a drainable porosity close to its specific yield (otherwise Raven warns that a rising water table
would release a different amount of water from the column than MODFLOW stores). Now add the columns:
\diff{t3_richards.diff}
\term{t3_run_richards.txt}
\term{t3_summary_rich.txt}
No new warning (\cmd{Raven\_errors.txt} holds the \V{t3:nwarnRich} warnings above). The run takes about \V{t3:secrichards}
seconds. Then the two other methods, each a one-line change of the file as the step before left it (run each into its own folder,
\cmd{out\_kinematic/} and \cmd{out\_lag/}):
\diff{t3_kinematic.diff}
\diff{t3_lag.diff}
And depth classes: in a cell on a slope the DEM points lie at different depths above the water table. With
\cmd{:SoilZoneVadoseClasses 3} each link gets up to three columns, one per depth class of its points (the change puts \cmd{RICHARDS} back and adds the
line; run into \cmd{out\_classes/}, about \V{t3:secclasses}~seconds):
\diff{t3_classes.diff}

\section{Compare the runs}
The tutorial's script \cmd{compare\_runs.py} prints the run means and draws the recharge reaching the water table, the
water-table depth of the first HRU given, and the streamflow:
\term{t3_compare.txt}
\fig{t3_compare.png}{The six runs: recharge reaching the water table (30-day means), water-table depth of the north slope
of the lower subbasin (HRU~7), and streamflow at the outlet.}
What the table says:
\begin{itemize}
\item \textbf{CONCEPTUAL to PHYSICAL.} In the conceptual run \V{t3:conceptual:leaves}~m\textsuperscript{3}/d leave the
  subsoil and all of it reaches MODFLOW; \V{t3:conceptual:seep}~m\textsuperscript{3}/d seep back out. With the soil zone,
  \V{t3:physical:leaves}~m\textsuperscript{3}/d leave the subsoil, because on the valley floor the soil now drains into an
  aquifer whose water table is in the soil, but \V{t3:physical:rejected}~m\textsuperscript{3} a day of it is rejected and
  \V{t3:physical:aqsoil}~m\textsuperscript{3}/d rises back from the aquifer into the soil: a two-way exchange replaces the
  one-way seepage drain (seepage to Raven falls to \V{t3:physical:seep}~m\textsuperscript{3}/d, from the lake HRU only). The
  water table under the valley floor (HRU~6) lies at \V{t3:physical:wt6}~m instead of \V{t3:conceptual:wt6}~m, and more
  water reaches the creek through the soil (\V{t3:physical:q} instead of \V{t3:conceptual:q}~m\textsuperscript{3}/s).
\item \textbf{The unsaturated columns} hold water: at the end of the run RICHARDS stores \V{t3:richards:uz}~million
  m\textsuperscript{3} above its equilibrium profile (\V{t3:uzmm}~mm over the \V{t3:slopearea}~km\textsuperscript{2} of
  slopes), LAG only \V{t3:lag:uz}~million m\textsuperscript{3} above the same profile, and KINEMATIC
  \V{t3:kinematic:uz}~million m\textsuperscript{3} (counted from the residual water content; manual, ``The unsaturated
  column''). The difference is physical: the columns start at rest, at equilibrium with a deep water table, and RICHARDS
  must wet that dry profile until it can carry the recharge. Over the three years it withholds \V{t3:withheld}\% of what
  enters the columns, so only \V{t3:richards:reaches}~m\textsuperscript{3}/d reach the water table, against
  \V{t3:kinematic:reaches} with KINEMATIC, and the water table under the slopes ends \V{t3:wtdrop}~m deeper than without
  columns. Most of that is spin-up (see the note below and \cmd{:SoilZoneVadoseInitialFlux}). Recharge arrives later and smoother (Figure~\ref{fig:t3vz}): the flow out of
  the columns follows the flow into them best \V{t3:richards:lag}~days later with RICHARDS, \V{t3:kinematic:lag} with
  KINEMATIC and \V{t3:lag:lag} with LAG.
\item \textbf{Depth classes} change little here (\V{t3:classes:reaches} against \V{t3:richards:reaches}
  m\textsuperscript{3}/d reaching the water table) and take longer to run. They matter where the unsaturated thickness
  within a cell varies a lot compared with its mean, as on hummocky ground over a shallow water table.
\end{itemize}
\fig{t3_vadose.png}{The unsaturated columns of the three methods: what enters them and what reaches the water table
(30-day means), and the water they store.\label{fig:t3vz}}

\begin{note}[A deep unsaturated zone needs a spin-up]
The columns start at rest (at equilibrium with the water table) and fill during the run: with RICHARDS the water table under
the slopes falls from \V{t3:physical:wt7}~m (no columns) to a mean of \V{t3:richards:wt7}~m, and the storage is still
growing after three years. Look at the column \cmd{unsaturated zone storage} of \cmd{GWBudget.csv}: as long as it trends,
the run is spinning up. Either run long enough and start the real run from its \cmd{solution.rvc} (Tutorial~5: the
columns are saved there), or start each column at its steady profile for a typical recharge with
\cmd{:SoilZoneVadoseInitialFlux} (mm/d; user guide, section~4.3c).
\end{note}

\section{Which to use}
\begin{itemize}
\item \cmd{CONCEPTUAL} when the water table stays below the soil, or for a first model.
\item \cmd{PHYSICAL} where the water table reaches the soil: valley floors, wetlands, toe slopes. It costs more per day.
\item \cmd{:SoilZoneVadose} where the water table lies metres below the soil and the timing of recharge matters:
      \cmd{RICHARDS} is the reference, \cmd{KINEMATIC} drops capillarity and cannot fail to converge, \cmd{LAG} is a cascade of
      linear stores with a travel time per metre.
\end{itemize}
The finished state, \cmd{\$TUT/t3\_soil\_zone/finished}, has the \cmd{.rvp} with all the parameters and the \cmd{.rvg} with
\cmd{RICHARDS} and three depth classes. The manual's section ``The soil zone and the water table'' gives the equations,
and the user guide's section~4.3c every \cmd{:SoilZone*} command.
